\chapter{Audiometry}

\section*{What Is This Test?}
Audiometry measures hearing sensitivity across different frequencies and loudness levels. Pure-tone audiometry is the standard behavioral hearing test and produces an audiogram showing air-conduction and, when needed, bone-conduction thresholds.

\section*{Alternative Names}
\begin{itemize}
  \item Pure-Tone Audiometry
  \item Hearing Test
  \item Audiogram
  \item Threshold Audiometry
\end{itemize}
\section*{Why This Test Is Ordered}
It evaluates hearing difficulty, tinnitus, occupational noise exposure, recurrent ear disease, suspected medication-related hearing loss, communication concerns, and hearing-aid or cochlear-implant needs. It also helps distinguish conductive from sensorineural hearing loss.

\section*{How This Test Is Performed}
The patient listens through headphones or insert earphones in a quiet booth and signals whenever a tone is heard. Thresholds are measured at selected frequencies for each ear. Bone-conduction testing uses a vibrator placed behind the ear when the type of hearing loss needs clarification. Speech recognition and tympanometry may be added.

\section*{Preparation, Risks, and Limitations}
No fasting or special preparation is needed. The test is painless and noninvasive. Background noise, poor attention, earwax, temporary middle-ear fluid, and inconsistent responses can affect the result. Results in young children may require age-appropriate play or conditioned testing.

\section*{Understanding Results}
Hearing thresholds are reported in decibels hearing level. The pattern and air-bone gap help classify hearing as normal, conductive, sensorineural, or mixed. Speech discrimination, asymmetry, and sudden loss may require prompt medical or specialist assessment.

\section*{References}
American Speech-Language-Hearing Association guidance on pure-tone and speech audiometry.

\clearpage
\section*{Understanding Results and Follow-up}
The laboratory report should identify the specimen, method, units, reference interval or decision limit, and whether the result is positive, negative, abnormal, or inconclusive. Reference intervals vary with age, sex, pregnancy, specimen type, assay, and laboratory; a result outside a range is not automatically a diagnosis, and a result within a range does not always exclude disease. Discuss unexpected results with the ordering clinician, who can decide whether repeat or confirmatory testing, treatment, or referral is appropriate. Do not change medicines or delay urgent care based on an isolated result.


\section*{Regulatory Status and Authorization Date}
This chapter describes a test category, not a specific commercial product. FDA, CDSCO, EU IVDR, and MHRA approval or registration dates therefore cannot be assigned without the assay manufacturer, product name, instrument, and intended use. For laboratory-developed tests, an individual agency approval date may not exist. See the Preface for the interpretation of regulatory status.

\clearpage
